\subsection{Dimension of an integral algebra}

LEMMA 2. --- Let \(\rho: A \to B\) be a ring homomorphism such that, for every prime ideal \(\mathfrak{p}\) of A, the \(\kappa(\mathfrak{p})\)-algebra \(B \otimes_{A} \kappa(\mathfrak{p})\) is integral (V, § 1, no. 1, def. 2). Let \(\mathfrak{q}\) and \(\mathfrak{q}'\) be two prime ideals of B such that \(\mathfrak{q} \subset \mathfrak{q}'\) and \(\mathfrak{q} \neq \mathfrak{q}'\). Then \(\rho^{-1}(\mathfrak{q}) \neq \rho^{-1}(\mathfrak{q}')\).

Indeed, if \(\mathfrak{q}\) and \(\mathfrak{q}'\) lie over the same prime ideal \(\mathfrak{p}\) of A, one has \(\dim(B \otimes_{A} \kappa(\mathfrak{p})) \geqslant 1\) by Lemma 1 of no. 1, which contradicts the fact that \(\dim(B \otimes_{A} \kappa(\mathfrak{p})) \leqslant 0\) (§ 1, no. 3, example 6).

THEOREM 1. --- Let \(\rho: A \to B\) be a ring homomorphism making \(B\)\(A\) an integral A-algebra.

\begin{enumerate}
\def\labelenumi{\alph{enumi})}
\item
  Let M be a finitely generated A-module. Then we have \(\dim_{\mathbf{B}}(\mathbf{M} \otimes_{\mathbf{A}} \mathbf{B}) \leqslant \dim_{\mathbf{A}}(\mathbf{M})\). In particular, we have \(\dim(\mathbf{B}) \leqslant \dim(\mathbf{A})\). If the map \({}^{a}\rho: \operatorname{Spec}(\mathbf{B}) \to \operatorname{Spec}(\mathbf{A})\) is surjective, for example (V, § 2, no. 1, theorem 1) if \(\rho\) is injective, we have \(\dim_{\mathbf{B}}(\mathbf{M} \otimes_{\mathbf{A}} \mathbf{B}) = \dim_{\mathbf{A}}(\mathbf{M})\), and in particular \(\dim(\mathbf{B}) = \dim(\mathbf{A})\).
\item
  Let  be an ideal of B and \(\mathfrak{a} = \rho^{-1}(\mathfrak{b})\) its inverse image in A. We have \(\operatorname{ht}(\mathfrak{b}) \leqslant \operatorname{ht}(\mathfrak{a})\) and \(\dim(\mathbf{B}/\mathfrak{b}) = \dim(\mathbf{A}/\mathfrak{a})\). If \({}^{a}\rho : \operatorname{Spec}(B) \to \operatorname{Spec}(A)\) is surjective, we have \(\operatorname{ht}(\mathfrak{a}\mathbf{B}) \leqslant \operatorname{ht}(\mathfrak{a})\) for every ideal  of A.
\item
  Suppose B is finite over A and let N be a finitely generated B-module. Then we have \(\dim_{\mathbf{B}}(\mathbf{N}) = \dim_{\mathbf{A}}(\mathbf{N})\) . In particular, we have \(\dim(\mathbf{B}) = \dim_{\mathbf{A}}(\mathbf{B})\) .
\end{enumerate}

Let us prove a). By prop. 5 of V, § 1, no. 1, the -algebra \(\kappa(\mathfrak{p})\) is integral for every prime ideal \(\otimes_{\mathbf{A}} \kappa(p)\)\(\mathfrak{p}\) of A. Let \(\mathfrak{q}_0 \subset \ldots \subset \mathfrak{q}_m\) be a chain of prime ideals of B; by Lemma 2, the ideals \(\mathfrak{p}_i = \rho^{-1}(\mathfrak{q}_i)\) are pairwise distinct, hence \(\mathfrak{p}_0 \subset \ldots \subset \mathfrak{p}_m\) is a chain of prime ideals of A, whence \(m \leqslant \dim(\mathbf{A})\) . One therefore has \(\dim(\mathbf{B}) \leqslant \dim(\mathbf{A})\) .

Suppose now that the map \({}^{a}\rho\) is surjective. Let \(\mathfrak{p}_0 \subset \ldots \subset \mathfrak{p}_n\) be a chain of prime ideals of A; there therefore exists a prime ideal \(\mathfrak{q}_0\) lying over \(\mathfrak{p}_0\). By cor. 2 to the first existence theorem (V, § 2, no. 1, th. 1), one can construct, by induction on \(n\), a chain \(\mathfrak{q}_0 \subset \ldots \subset \mathfrak{q}_n\) of prime ideals of B such that \(\mathfrak{q}_i\) lies over \(\mathfrak{p}_i\) for \(0 \leqslant i \leqslant n\). One therefore has \(n \leqslant \dim(\mathbf{B})\) and consequently \(\dim(\mathbf{A}) \leqslant \dim(\mathbf{B})\).

This proves \(a\) ) in the case where \(\mathbf{M} = \mathbf{A}\) . In the general case, let \(\alpha\) denote the annihilator of \(\mathbf{M}\), so that the support of \(\mathbf{M}\) is identified with , and one has \(\operatorname{Spec}(\mathbf{A}/\alpha)\)\(\dim_{\mathbf{A}}(\mathbf{M}) = \dim(\mathbf{A}/\alpha)\). By II, § 4, no. 4, prop. 19, the support of \(\mathbf{M} \otimes_{\mathbf{A}} \mathbf{B}\) is the inverse image under \(^{\alpha}\rho\) of the support of \(\mathbf{M}\), hence is identified with , and we have \(\operatorname{Spec}(\mathbf{B}/\rho(\alpha)\mathbf{B})\). It remains to note that the homomorphism \(\dim_{\mathbf{B}}(\mathbf{M} \otimes_{\mathbf{A}} \mathbf{B}) = \dim(\mathbf{B}/\rho(\alpha)\mathbf{B})\)\(\rho': A/\alpha \to B/\rho(\alpha)\) deduced from \(\rho\) makes  an integral \(B/\rho(\alpha)\)-algebra, and that \((A/\alpha)\) is surjective when \(^{\alpha}\rho'\)\(^{\alpha}\rho\) is.

We prove \(b\) ). By prop. 7 of § 1, no. 3, it suffices to prove that  for every prime ideal  of A containing ; let  be such an ideal. By V, § 2, no. 1, cor. 2 to th. 1, there exists a prime ideal \_ of B lying over  and containing , and one has \_ by prop. 7 of § 1, no. 3.

Now \(\mathbf{B}_{\mathfrak{q}}\) is identified with a ring of fractions of the integral \(\mathbf{A}_{\mathfrak{p}}\) -algebra \(\mathbf{B} \otimes_{\mathbf{A}} \mathbf{A}_{\mathfrak{p}}\) , whence

\[
\dim (\mathrm{B} _ {\mathfrak {q}}) \leqslant \dim (\mathrm{B} \otimes_ {\mathrm{A}} \mathrm{A} _ {\mathfrak {p}}) \leqslant \dim (\mathrm{A} _ {\mathfrak {p}})
\]

by prop. 6 of § 1, no. 3 and assertion a) above. We have thus proved the inequality . Moreover, the homomorphism of  into  deduced from  is injective and makes  an integral -algebra; hence one has  by a). Suppose  is surjective, and let  be an ideal of A and  a prime ideal of A containing . By hypothesis there exists a prime ideal  of B lying over . We have , whence  by what precedes. Passing to the infimum, we obtain .

Finally,  follows from  applied to the annihilator  of N.

THEOREM 2. --- Let A be an integrally closed ring, and B a ring containing A, integral over A. Suppose that B is a torsion-free A-module. For every ideal  of A, one has . Let  be an ideal of B and ; then one has .

Let \(\rho\) be the canonical map from A to B. Let \(\mathfrak{a}\) be an ideal of A. If \(\mathfrak{a} = \mathbf{A}\), the first equality is clear. Suppose \(\mathfrak{a} \neq \mathbf{A}\). Since \(\rho\) is injective, \({}^{a}\rho\) is surjective (V, § 2, no. 1, th. 1). Consequently \(\mathfrak{aB} \neq \mathbf{B}\). Now let \(\mathfrak{q}\) be a prime ideal of B containing \(\mathfrak{aB}\). Put \(\mathfrak{p} = \mathfrak{q} \cap \mathbf{A}\). We have \(\mathfrak{a} \subset \mathfrak{p}\), whence . Let \(\mathfrak{p}_0 \subset \ldots \subset \mathfrak{p}_n\) be a chain of prime ideals of A with \(\mathfrak{p}_n = \mathfrak{p}\). By the second existence theorem (V, § 2, no. 4, th. 3), one constructs by induction a chain \(\mathfrak{q}_0 \subset \ldots \subset \mathfrak{q}_n\) of prime ideals of B such that \(\mathfrak{q}_n = \mathfrak{q}\) and \(\mathfrak{q}_i\) lies over \(\mathfrak{p}_i\) for \(0 \leqslant i \leqslant n\). We have , whence . Passing to the infimum, one obtains  (§ 1, no. 3, prop. 7). The inequality  follows from th. 1, whence the first equality. Let \(n \leqslant\) be an ideal of B. Put \(\mathfrak{a} = \rho^{-1}(\mathfrak{b})\). We have , whence . The inequality  follows from th. 1, whence the theorem.

Remark. --- Let A be an integral domain and B a ring containing A, integral over A. Let  be a prime ideal of A such that the integral closure of  is a local ring. One can show that, for every prime ideal  of B lying over , one has \(A_{p}\) (p. 85, exercise 9) when B is an integral domain.

